\section{Catálogo de relaciones}\label{sec:relaciones}

En este catálogo defino las relaciones que uso en las expresiones.
El índice que sigue al signo K03 selecciona la relación. Sus argumentos
ocupan las posiciones indicadas y se separan con una barra vertical.
Los índices de relaciones y los de referentes pertenecen a catálogos diferentes.

\subsection{Relación 1: precede temporalmente}

\textbf{Posiciones:} 1. evento anterior; 2. evento posterior.

En el marco temporal declarado, el primer evento termina antes de que comience el segundo. Para eventos puntuales se comparan sus instantes. No afirma causalidad.

\subsection{Relación 2: lee un dato}

\textbf{Posiciones:} 1. agente; 2. dirección; 3. dato.

En el evento descrito, el agente lee el dato indicado de la dirección indicada. Se usa como descripción abstracta; no fija modelo de memoria, anchura del dato ni concurrencia.

\subsection{Relación 3: forma parte de}

\textbf{Posiciones:} 1. parte; 2. todo.

En el modelo declarado, el primer objeto es una parte propia del segundo. Es una relación de parte a todo, distinta de pertenecer a un conjunto.

\subsection{Relación 4: tiene mayor frecuencia}

\textbf{Posiciones:} 1. primer tono; 2. segundo tono.

La frecuencia fundamental del primer tono es mayor que la del segundo, en el intervalo y las condiciones que se declaren. No afirma mayor volumen ni mayor duración.

\subsection{Relación 5: pertenece a un conjunto}
\textbf{Posiciones:} 1. elemento; 2. conjunto o extensión de un concepto.
Es verdadera si el primer término es miembro del segundo en el contexto declarado.
No expresa parte propia. Una lista no es un segundo argumento válido en este fragmento.
Ejemplo: T08; \(A\in\{A,B\}\).

\subsection{Relación 6: influencia causal en un modelo}
\textbf{Posiciones:} 1. variable X; 2. variable Y; 3. modelo causal M.
El dominio admitido son modelos estructurales deterministas acíclicos, con variables
endógenas distintas X e Y. M declara variables, dominios, ecuaciones y contextos
exógenos admisibles. R6 es verdadera cuando existe un mismo contexto \(u\) y dos
valores admisibles \(x_0,x_1\) tales que
\[
 Y_{\operatorname{do}(X=x_0)}(u)\ne
 Y_{\operatorname{do}(X=x_1)}(u).
\]
La intervención sustituye la ecuación de X por el valor fijado y conserva las
demás ecuaciones. Es una influencia causal posible en el modelo, no una
afirmación de causa efectiva de un evento histórico. No se infiere únicamente de
correlación o precedencia temporal. Ejemplo completamente definido: T10.

\subsection{Relación 7: le gusta un libro}
\textbf{Posiciones:} 1. persona; 2. libro.
La persona satisface el criterio de gusto respecto de ese libro en el estudio
declarado, por ejemplo una respuesta afirmativa a una pregunta fija en una fecha.
No significa que sea su favorito, que lo haya comprado o que lo posea.
Ejemplos: T04 y T05, con C1 y C2.

\subsection{Relación 8: satisface una propiedad}
\textbf{Posiciones:} 1. objeto; 2. propiedad indexada.
El objeto cumple el criterio de la propiedad en el contexto fijado para toda la
expresión. Las condiciones temporales o experimentales exigidas por la propiedad
deben declararse. No indica intensidad ni probabilidad. Ejemplo: T16, con P1 y P2.

\subsection{Relación 9: está al norte en un marco}
\textbf{Posiciones:} 1. lugar A; 2. lugar B; 3. marco geográfico F.
Los lugares tienen coordenadas válidas en F, con datum y época declarados. R9
significa que la latitud de A es estrictamente mayor que la de B, usando latitud
positiva hacia el norte. No afirma distancia ni vecindad. Ejemplo: T18.
